\documentclass[a4paper, 11pt]{article}
\usepackage[utf8]{inputenc}
\usepackage[T1]{fontenc}
\usepackage[italian]{babel}
\usepackage{geometry}
\usepackage{graphicx}
\usepackage{tabularx}
\usepackage{fancyhdr}
\usepackage{hyperref}

% Margini
\geometry{top=3cm, bottom=3cm, left=2.5cm, right=2.5cm}

% Intestazione e piè di pagina
\pagestyle{fancy}
\fancyhf{}
\fancyhead[L]{\small \textit{0x06tudio} -- Verbale di Riunione}
\fancyhead[R]{\small Anno Accademico 2026/2027}
\fancyfoot[C]{\thepage}
\renewcommand{\headrulewidth}{0.4pt}

\begin{document}
\thispagestyle{plain}

\newcommand{\issue}[1]{\href{https://github.com/0x06tudio/docs/issues/#1}{\textit{#1}}}

\begin{minipage}{0.6\textwidth}
    \vspace{-2cm}
    {\huge \textbf{0x06tudio}} \\[0.5em]
    {\large Progetto didattico di Ingegneria del Software} \\[0.3em]
    {\normalsize Università degli Studi di Padova}
\end{minipage}
\begin{minipage}{0.4\textwidth}
    \raggedleft
    \vspace{0pt}
    \includegraphics[width=6cm]{../../media/logo.png}
    \vspace{1cm}
    \small
    \href{mailto:0x06tudio@gmail.com}{0x06tudio@gmail.com}
\end{minipage}

\vspace{0.5cm}

\begin{center}
    {\LARGE \textbf{Verbale Riunione}}
\end{center}

\vspace{0.5cm}


\begin{tabularx}{\textwidth}{l X}
    \textbf{Versione:} & 0.1.0 \\
    \textbf{Stato:} & Non approvato \\
    \textbf{Data:} & 2026/11/13 \\
    \textbf{Uso:} & Esterno \\
    \textbf{Partecipanti (Presenti):} & Nicolò Beltrame, Pietro Brocadello, Alberto Cavalletto, Catalin Chirilov, Davide Colabove, Andrea Maggio \\
    \textbf{Destinatari:} & Zucchetti SpA, Proff Vardanega, Cardin \\
    \textbf{Modalità / Luogo:} & Online / Google Meet \\
    \textbf{Ora inizio:} & 15:00 \\
    \textbf{Ora fine:} & 16:00 \\
    \textbf{Durata totale:} & 1 ora \\
    \textbf{Partecipanti esterni:} & Gregorio Piccoli (Zucchetti SpA)\\
\end{tabularx}
\section*{Ordine del giorno}
\begin{enumerate}
    \item Discussione con il referente esterno riguardo le tecnologie da esplorare e studiare autonomamente prima dell'inizio dei lavori\vspace{-0.2cm}
    \item Chiarimenti sul perimetro del progetto \vspace{-0.2cm}
    \item Proposta di un piano per il PoC \vspace{-0.2cm}
\end{enumerate}
\vspace{-0.5cm}

\section*{Registro delle modifiche}
\begin{tabularx}{\textwidth}{|c|c|X|c|c|}
    \hline
    \textbf{Versione} & \textbf{Data} & \textbf{Descrizione} & \textbf{Autore} & \textbf{Verificatore} \\ \hline
    0.1.0 & 2026/11/14 & Stesura del documento & Andrea Maggio & / \\ \hline
\end{tabularx}

\newpage
\tableofcontents
\newpage


\section{Resoconto della discussione}

\subsection{Punto 1: Discussione sulle tecnologie}
Il gruppo ha avanzato varie proposte su tecnologie da esaminare, dettagliate nelle sub-issue della issue \issue{28} della repository del gruppo dedicata alla documentazione. Il referente esterno ha approvato le proposte del gruppo, e ha offerto alcuni spunti di riflessione riguardo lo Spec-Driven development:
\begin{itemize}
    \item \href{https://ranthebuilder.cloud/blog/i-tested-three-spec-driven-ai-tools-here-s-my-honest-take}{\texttt{ranthebuilder.cloud/blog/}} \vspace{-0.2cm}
    \item Una relazione di un precedente stagista, che il referente invierà via mail al recapito del gruppo\vspace{-0.2cm}
\end{itemize}
\subsection{Punto 2: Chiarimenti sulle richieste per il progetto}
Il referente esterno ha mostrato al gruppo il suo attuale flusso di lavoro che usufruisce di openCode con openSpec per mostrare chiaramente le sue limitazioni e quali obiettivi vorrebbe dall'MVP. In particolare il referente ha posto particolare attenzione alla mancanza di rigorosità e formalismi nella pratica del \textit{vibe coding}: l'obiettivo principale{\renewcommand{\thefootnote}{*} \footnote{il progetto include anche altre richieste, descritte nel documento di presentazione del capitolato. Il chiarimento mostrato in seguito è da intendersi come spiegazione "ad alto livello" e non esaustiva.}} è quello di sostituire il prompt con la lettura di un diagramma dei casi d'uso, e anche verificare che il prodotto finale effettivamente soddisfi tali casi d'uso. In tal modo è possibile evitare l'ambiguità del linguaggio naturale.

Inoltre il referente ha fornito chiarimenti sul funzionamento dei \textit{code agents}, dando alcune definizioni:
\begin{itemize}
    \item \textbf{\textit{skill:}} una skill è un file \texttt{.md} scritto in linguaggio naturale che spiega all'LLM la procedura, le best practices e l'ordine in cui effettuare una specifica richiesta. Il risultato è certamente non deterministico: sono delle aggiunte al contesto dell'LLM, che dunque le interpreterà arbitrariamente. Importante notare che le skill non vengono iniettate per intero in ogni prompt: ne viene iniettata una breve descrizione, così il modello potrà decidere autonomamente, se ritiene necessario, di leggerla.   \vspace{-0.2cm}
    \item \textbf{\textit{tool:}} un tool è un file strutturato (nell'esempio con openSpec un \texttt{JSON}) che fornisce una descrizione di uno strumento specifico, e i parametri da passargli (ad esempio \texttt{run\_bash(comando)}). Dunque il tool è deterministico nella sua esecuzione, ma l'LLM decide quando e con quali parametri usarlo. I tool vengono sempre iniettati nel prompt, perché il modello deve sempre sapere cosa può fare e come farlo. \vspace{-0.2cm}
    \item \textbf{\textit{comando:}} nel contesto di openSpec, i comandi con sintassi \texttt{/opsx-nome} sono tecnicamente delle skill, ma attivate manualmente su necessità: ad esempio si usa \texttt{/opsx-propose} per obbligare il modello a seguire un flusso di lavoro specifico per il quale genera una \textit{documentazione di pianificazione} prima di agire sul codice. \vspace{-0.2cm}
\end{itemize}

\subsection{Punto 3: Piano per il PoC}
A causa dell'esaurimento della quantità di tempo disponibile al referente per la riunione, la proposta di un piano di azione per la realizzazione del PoC è stata ridotta ad una e-mail, di cui il gruppo attende risposta.
\newpage


\section{Approvazione esterna}
Il presente verbale è stato condiviso con il referente aziendale in data YYYY/MM/DD ed è da ritenersi approvato e conforme a quanto discusso durante il colloquio. 


\end{document}